%------------------------------------------------------------------------------%

\usepackage{integracao_e_series}

\title{Integrais Impróprias}

%------------------------------------------------------------------------------%
\begin{document}

\addtitle

%------------------------------------------------------------------------------%
\section{Integrais Impróprias}

%------------------------------------------------------------------------------%
\begin{frame}{Integrais Impróprias}
\begin{enumerate}[<+->]
  \setlength\itemsep{1em}
  \item Intervalo de integração infinito
  \item Descontinuidade no extremo do intervalo
\end{enumerate}
\end{frame}

%------------------------------------------------------------------------------%
\begin{frame}{Integrais Impróprias}
\begin{enumerate}[<+->]
  \setlength\itemsep{1em}
  \item Envolve o calculo de um limite
  \item Se o limite existe a integral é \emph{convergente}
  \item Se algum limite não existir a integral é \emph{divergente}
\end{enumerate}
\end{frame}

%------------------------------------------------------------------------------%
\section{Integrais Impróprias -- Tipo 1}

%------------------------------------------------------------------------------%
\begin{frame}{Integrais Impróprias -- Tipo 1}
  $f$ função contínua em intervalo ilimitado,
  \emph{\([a, \infty)\)},
  \emph{\((-\infty, b]\)},
  \emph{\((-\infty, \infty)\)}
\end{frame}

%------------------------------------------------------------------------------%
\begin{frame}{Integrais Impróprias -- Tipo 1}
  $f$ é função contínua em \([a, \infty)\)
  \large
  \[
    \int_a^{{\color<2>{blue} \infty}} f(x) \, dx
    \;=\;  {{\color<2>{blue} \lim_{t\to\infty}}}
    \int_a^{{\color<2>{blue} t}} f(x)  \, dx
  \]
\end{frame}

%------------------------------------------------------------------------------%
\begin{frame}{Integrais Impróprias -- Tipo 1}
  $f$ é função contínua em \((-\infty, b]\)
  \large
  \[
    \int_{{\color<2>{blue} -\infty}}^b f(x) \, dx
    \;=\; {\color<2>{blue} \lim_{t\to-\infty}}
    \int_{{\color<2>{blue} t}}^{b} f(x) \,dx
  \]
\end{frame}


%------------------------------------------------------------------------------%
\begin{frame}{Integrais Impróprias -- Tipo 1}

  $f$ é função contínua em \((-\infty, \infty)\)

  {\large
  \[
    \int_{{\color<3>{blue}  -\infty}}
        ^{{\color<4>{magenta}\infty}} f(x) \, dx
    \;=\; {\color<3>{blue}   \lim_{r\to-\infty}}
    \int_{{\color<3>{blue}{r}}}
        ^{{\color<2>{red} c}} f(x) \, dx
    \;+\; {\color<4>{magenta}\lim_{s \to \infty}}
    \int_{{\color<2>{red} c}}
        ^{{\color<4>{magenta}s}} f(x) \, dx
  \]}
  para algum {\color<2>{red} $c$} $\in\R$

\end{frame}

%------------------------------------------------------------------------------%
\section{Integrais Impróprias -- Tipo 2}

%------------------------------------------------------------------------------%
\begin{frame}{Integrais Impróprias -- Tipo 1}
  $f$ função \emph{contínua no interior} do intervalo,
  mas tende ao \emph{infinito} em algum ponto
\end{frame}

%------------------------------------------------------------------------------%
\begin{frame}{Integrais Impróprias -- Tipo 2}
  $f$ é contínua em {\color{blue} $[a, b)$} e descontínua em $b$,
  {
  \large
  \[
    \int_{a}^{{\color<2>{blue} b}} f(x)\,dx
    \;=\;     {\color<2>{blue} \lim_{t\to b^{-}}}
    \int_{a}^{{\color<2>{blue} t}} f(x)\,dx
  \]}
  Atenção ao limite lateral
\end{frame}

%------------------------------------------------------------------------------%
\begin{frame}{Integrais Impróprias -- Tipo 2}
  $f$ é contínua em {\color{blue} $(a, b]$} e descontínua em $a$,
  {
  \large
  \[
    \int_{{\color<2>{blue} a}}^b f(x)\,dx
    \;=\; {\color<2>{blue} \lim_{t\to a^+}}
    \int_{{\color<2>{blue} t}}^b f(x)\,dx
  \]}
  Atenção ao limite lateral
\end{frame}

%------------------------------------------------------------------------------%
\begin{frame}{Integrais Impróprias -- Tipo 2}
  $f$ é uma função definida em $[a, b]$, contínua exceto em um ponto
  {\color{blue} $c \in (a, b)$}
  {
  \large
  \begin{align*}
      \int_a^b f(x) \, dx
    & \;=\; \int^{{\color<2>{blue} c}}_a f(x) \, dx
      \;+\; \int_{{\color<2>{blue} c}}^b f(x) \, dx \\[3mm]
    & \;=\; {\color<2>{blue} \lim_{r\to c^-}} \int^{{\color<2>{blue} r}}_a f(x) \, dx
      \;+\; {\color<2>{blue} \lim_{s\to c^+}} \int_{{\color<2>{blue} s}}^b f(x) \, dx
  \end{align*}

  }
\end{frame}

%------------------------------------------------------------------------------%
\section{Exemplos}

%------------------------------------------------------------------------------%
\addtocounter{exemplo}{1}
\begin{frame}{Exemplo \theexemplo}
  Calcular a integral \(\int_{-\infty}^0 e^x\, dx\)
\end{frame}

\begin{frame}{Exemplo \theexemplo}
\begin{align*}
  \onslide<+->{\int_{{\color<2>{blue}-\infty}}^0 e^x\, dx
                 & = {\color<2>{blue}\lim_{t\to-\infty}}
       \left(\int_{{\color<2>{blue} t}}^{0} e^x\, dx \;\right)
    & & \text{definição integral imprópria}                       \\}
  \onslide<+->{
    & = \lim_{t\to-\infty}
    \left(e^x \,\evalat{t}{0} \;\right) \\[2mm]}
  \onslide<+->{
    & = \lim_{t\to-\infty} \left( e^0 - e^t \right) \\[2mm]}
  \onslide<+->{
    & = \lim_{t\to-\infty} 1 - \lim_{t\to-\infty}e^t \\[2mm]}
  \onslide<+->{& = 1}
\end{align*}
\end{frame}

%------------------------------------------------------------------------------%
\addtocounter{exemplo}{1}
\begin{frame}{Exemplo \theexemplo}
  Determine se a integral
  \[
    \int_{-1}^{0} \frac{1}{\,x^{\sfrac{2}{3}}\,}\;dx
  \]
  é convergente ou divergente
\end{frame}

\begin{frame}{Exemplo \theexemplo}
  \onslide<+->{\(f(x) = \frac{1}{\,x^{\sfrac{2}{3}}} = x^{\sfrac{-2}{3}}\)
    é descontínua em $x=0$}
  \begin{align*}
    \onslide<+->{\int_{-1}^{0} \frac{1}{\,x^{\sfrac{2}{3}}\,}\;dx}
    \onslide<+->{& = \lim_{t \to 0^{-}} \int_{-1}^{t} x^{\sfrac{-2}{3}}\;dx  \\[2mm]}
    \onslide<+->{& = \lim_{t \to 0^{-}} \left(\frac{x^{\sfrac{-2}{3}+1}}{\sfrac{-2}{3}+1} \evalat{-1}{t}\right) \\[2mm]}
    \onslide<+->{& = \lim_{t \to 0^{-}} \left(3x^{\sfrac{1}{3}} \evalat{-1}{t}\right) \\[2mm]}
    \onslide<+->{& = \lim_{t \to 0^{-}} 3 \left[ t^{\sfrac{1}{3}} - (-1)^{\sfrac{1}{3}} \right] \\[2mm]}
    \onslide<+->{& = 3}
    \onslide<+->{& \text{Portanto, a integral imprópria converge}}
  \end{align*}
\end{frame}

%------------------------------------------------------------------------------%
\addtocounter{exemplo}{1}
\begin{frame}{Exemplo \theexemplo}
  Calcule
  \[
    \int_{0}^3 \frac{1}{x-1} \, dx
  \]
  se possível
\end{frame}

\begin{frame}{Exemplo \theexemplo}
  A função \(f(x) = \frac{1}{x-1}\) é descontínua em $x=1\in[0, 3]$

  \pause
  \vspace{\baselineskip}
  Precisamos dividir o intervalo e duas partes \([0, 1]\) e \([1, 3]\)
  \pause
  \[
          \int_0^3 \frac{1}{x-1} \,dx
    \;=\; \int_0^1 \frac{1}{x-1} \,dx
    \;+\; \int_1^3 \frac{1}{x-1} \,dx
  \]

  \pause
  Agora a descontinuidade está nos extremos dos intervalos
  e podemos usar integrais impróprias
\end{frame}

\begin{frame}{Exemplo \theexemplo}
\begin{align*}
  \onslide<+->{\int_{0}^1 \frac{1}{x-1} dx}
  \onslide<+->{& = \lim_{t \to 1^{-}}\int_{0}^t \frac{1}{x-1} dx  \\}
  \onslide<+->{& = \lim_{t \to 1^{-}} \ln\abs{x-1} \evalat{0}{t}  \\}
  \onslide<+->{& = \lim_{t \to 1^{-}}\big(\ln\abs{t-1} - \ln\abs{-1} \big) \\}
  \onslide<+->{& = \lim_{t \to 1^{-}} \ln(1-t) }
  \onslide<+->{& & 1-t \to 0^+ \text{ quando } t \to 1^{-} \\}
  \onslide<+->{& =-\infty}
\end{align*}
\onslide<+->{Portanto a integral diverve e não precisamos calcular a segunda parte}
\end{frame}

%------------------------------------------------------------------------------%
\section{Lista Mínima}

%------------------------------------------------------------------------------%
\begin{frame}{Lista Mínima}
  Estudar a Seção 2.6 da Apostila

  Fazer os exercícios 1, 2 e 3

  \vfill
  Atenção: A prova é baseada no livro, não nas apresentações
\end{frame}

%------------------------------------------------------------------------------%
\end{document}
%------------------------------------------------------------------------------%
